%% ch1_introduction.tex

\chapter{INTRODUCTION}

Fifth Generation (5G) mobile networks introduce network slicing as a defining architectural capability, enabling multiple logical network instances---each with distinct performance characteristics---to coexist on shared physical infrastructure. Network slicing permits concurrent support of enhanced Mobile Broadband (eMBB), Ultra-Reliable Low-Latency Communication (URLLC), and massive Machine-Type Communication (mMTC) services, each with fundamentally different QoS requirements. eMBB demands high throughput and tolerates moderate latency; URLLC requires end-to-end latencies below 20\,ms with very high reliability; and mMTC targets energy-efficient transmission for large-scale IoT device populations.

Realising these guarantees in practice is challenging. At the user plane, traffic from different slices may contend for shared buffering and forwarding resources during congestion. Static resource allocations fail to track dynamic traffic demands, while reactive threshold-based controllers lack the ability to reason about the root cause of degradation or to anticipate violations before they occur.

Autonomous orchestration, enabled by Large Language Models (LLMs) and multi-agent frameworks, offers a qualitatively different control paradigm. Rather than pattern-matching against hardcoded thresholds, an LLM-based planning agent performs causal reasoning: identifying which network entity is responsible for congestion, validating whether the proposed control action addresses that cause, and consulting prior operational history stored in an agent memory to refine its decisions. This project implements, deploys, and evaluates such a system in a cloud-native 5G testbed.

\section{Motivation}

The evolution from rule-based to autonomous orchestration is driven by three limitations observed in practical 5G slice management systems.

\textbf{First}, rule-based controllers are structurally blind to causation. A threshold-based system that detects elevated URLLC RTT will throttle eMBB bandwidth regardless of whether eMBB traffic is actually active. When eMBB load is low, this action is unnecessary and harmful. An agent capable of reasoning about relative load---quantified by the \texttt{embb\_load\_fraction} metric introduced in this work---can correctly distinguish idle from bursting eMBB and decline to act when throttling is ineffective.

\textbf{Second}, reactive control cannot prevent violations; it can only respond to them. A URLLC SLA breach lasting even 200\,ms may disrupt a latency-critical industrial automation session. Pre-emptive action, triggered by trend analysis and predictive causal reasoning, is only possible when the controller possesses a causal model of the network state rather than a simple threshold table.

\textbf{Third}, accumulated operational experience is wasted in memoryless systems. An agent that stores the outcome of prior actions---including whether throttling eMBB at a given load level subsequently resolved the RTT violation---can converge to more accurate control policies over time without requiring offline training.

\section{Literature Survey}

\textbf{Bandara et al.} propose an AI-native framework for 6G network slice orchestration, monitoring, and trading using multiple autonomous AI agents and Large Language Models (LLMs)~\cite{bandara2026}. The architecture consists of layered components that perform intent interpretation, slice deployment, SLA monitoring, policy enforcement, and resource allocation. The framework integrates the Model Context Protocol (MCP) to translate natural language requests into secure orchestration actions and was evaluated on an Open5GS, Ericsson RAN, and Kubernetes testbed.

\textbf{Grings et al.} propose a full dynamic orchestration framework for 5G core network slicing over a cloud-native Kubernetes platform~\cite{grings2022}. Unlike existing approaches that only support partial orchestration through resource scaling, the proposed solution introduces a Kubernetes-integrated controller capable of dynamically reconfiguring 5G core functions, including NSSF, AMF, SMF, and UPF, during runtime without terminating active network slices, reducing reconfiguration requests by approximately 47.5\%.

\textbf{Novanana et al.} propose a Kubernetes-based MANO framework for provisioning coexisting eMBB and URLLC services through 5G network slicing using open-source technologies~\cite{novanana2024}. The authors integrate free5GC, UERANSIM, OpenStack, and kube5gnfvo to evaluate separate UPF vs. dedicated SMF/UPF configurations, showing that dedicated SMF/UPF pairs deliver better performance, achieving a stable throughput of 20.2--22.1\,Gbps and improved congestion control.

\textbf{Reddy} proposes a Deep Reinforcement Learning (DRL)-based Kubernetes autoscaling framework for managing high-throughput 5G network functions and slices using a Deep Q-Learning (DQN) agent~\cite{reddy2025}. The system dynamically scales Kubernetes pods based on real-time network conditions. Experimental results show 71.7\% faster slice deployment, 35\% higher resource utilization, 22\% lower latency, 40\% faster processing times, and up to 96.7\% spectrum efficiency compared to conventional methods.

\textbf{Dandoush et al.} propose an LLM-powered network slicing management and orchestration framework that integrates Large Language Models (LLMs) and multi-agent systems with ETSI MANO and 3GPP slicing architectures~\cite{dandoush2024}. The conceptual framework enables LLM agents to translate natural language intent into technical slice requirements, automate slice design and deployment, coordinate multi-domain resources, and manage performance without needing hardcoded rules.

\textbf{Sulaiman et al.} propose MicroOpt, a model-driven framework for dynamic resource optimization in 5G network slicing that uses a differentiable deep neural network (DNN) and gradient-based optimization to allocate resources~\cite{sulaiman2024}. The system predicts traffic, models QoS, and uses a primal-dual optimization algorithm to dynamically scale CPU and bandwidth. Deployed on an open-source testbed with srsRAN, Open5GS, Kubernetes, and Prometheus, it achieved up to 21.9\% lower resource consumption.

\textbf{Saha et al.} propose MonArch, a scalable monitoring architecture for cloud-native 5G deployments that integrates Prometheus, Elasticsearch, Apache Kafka, Logstash, and Kubernetes to collect, correlate, and analyze monitoring data across slices~\cite{saha2022}. Tested on a Kubernetes-based 5G testbed using Free5GC and UERANSIM, MonArch demonstrated minimal data-ingestion overhead and efficient scaling using a 5-second monitoring interval.

\textbf{Tran et al.} propose PCLANSA, a machine-learning-based closed-loop service assurance framework for 5G and B5G network slicing~\cite{tran2025}. The framework uses LSTM-based traffic prediction to forecast future demand and automatically perform VNF scaling and link-capacity adjustments, achieving resource savings of 54.85\% for eMBB, 57.1\% for uRLLC, 50.87\% for mMTC, and 23.63\% for VoIP slices in an OMNeT++ simulation.

\section{Problem Statement}

The objective of this project is to design, implement, and evaluate an autonomous agent-based orchestration framework for 5G network slicing that: (i) enforces user-plane QoS isolation across eMBB, URLLC, and mMTC slices through dynamic Linux \texttt{tc} shaping on Kubernetes-hosted UPFs; (ii) employs an LLM-assisted multi-agent architecture for monitoring, planning, and execution using structured causal reasoning; and (iii) evaluates its ability to improve QoS protection and reduce unnecessary control actions compared to a rule-based baseline through a pre-defined three-scenario behavioral audit.

\section{Objectives and Scope of the Project}

\subsection{Objectives}

\begin{enumerate}[leftmargin=*]
  \item To design and deploy a cloud-native 5G network slicing testbed using Open5GS SA Core, UERANSIM, Kubernetes, and per-slice UPF instances with user-plane isolation enforced through Linux \texttt{tc} HTB traffic shaping.

  \item To implement a LangGraph-based multi-agent orchestration framework consisting of Monitoring, Planning, Execution, and Memory agents operating within an asynchronous control loop.

  \item To develop and validate a structured reasoning framework incorporating root-cause assessment, lever validation, and workload-awareness metrics such as \texttt{embb\_load\_fraction} for informed orchestration decisions.

  \item To conduct a three-scenario behavioral audit comparing the proposed agent-based orchestrator with a rule-based baseline in terms of pre-emptive throttling, wrong-lever avoidance, and memory-assisted restoration.

  \item To evaluate the effectiveness of the proposed framework using metrics including URLLC RTT SLA compliance, eMBB throughput degradation, mMTC packet delivery ratio, recovery time, and unnecessary throttle rate.
\end{enumerate}

\subsection{Scope of the Project}

The project focuses on a single-domain 5G slice orchestration system operating at the user-plane level through dynamic bandwidth shaping. The orchestration intelligence is provided by a locally hosted LLM to preserve data sovereignty and eliminate cloud API dependencies. Extensions to multi-domain slicing, RL-based policy refinement, and O-RAN xApp integration are identified as directions for future work.

%% =============================================================
\chapter{REQUIREMENT ANALYSIS}

\section{Functional Requirements}

\begin{itemize}[leftmargin=*]
  \item The system shall support the simultaneous operation of three 5G network slices (eMBB, URLLC, mMTC) on shared Kubernetes infrastructure with enforceable user-plane isolation.
  \item The system shall establish concurrent PDU sessions per UE across all three slices using UERANSIM.
  \item The system shall enforce per-slice bandwidth caps using Linux \texttt{tc} HTB shaping applied to per-slice UPF GTP-U interfaces.
  \item The system shall continuously monitor per-slice QoS metrics with a 3-second collection interval via a dedicated Monitoring Agent.
  \item The system shall compute \texttt{embb\_load\_fraction} from a rolling window of packet-rate observations.
  \item The system shall invoke an LLM-backed Planning Agent producing a Chain-of-Thought trace (root cause, lever validity, action, confidence) for every cycle.
  \item The system shall detect internal reasoning contradictions and downgrade action confidence automatically.
  \item The system shall persist decision history in an Agent Memory and surface prior decisions in the LLM prompt context.
  \item The system shall execute validated actions via an Execution Agent applying \texttt{tc} commands via SSH.
  \item The system shall expose orchestrator state metrics via Prometheus on port 9200 and visualise them through Grafana.
\end{itemize}

\section{Non-Functional Requirements}

\begin{itemize}[leftmargin=*]
  \item The monitoring thread shall operate at 3-second interval, independent of LLM inference latency.
  \item LLM inference shall complete within 60\,s; a rule-based fallback activates upon timeout.
  \item \texttt{tc} shaping commands shall take effect within 200\,ms of issuance.
  \item The Agent Memory shall retain the 10 most recent decision entries.
  \item All decisions shall be logged with full Chain-of-Thought traces to support post-hoc audit.
  \item The system shall support dry-run mode for experimental evaluation without applying \texttt{tc} commands.
\end{itemize}

\section{Hardware Requirements}

\begin{table}[H]
\centering
\caption{Hardware Configuration of the Testbed}
\label{tab:hardware}
\begin{tabular}{llll}
\toprule
\textbf{Node} & \textbf{Role} & \textbf{RAM} & \textbf{Storage} \\
\midrule
kubemaster (192.168.49.174) & K8s Control Plane / Orchestrator & 16\,GB & 256\,GB \\
kube       (192.168.49.173) & K8s Worker 1 / Primary UPF pods  & 16\,GB & 256\,GB \\
kube2      (192.168.49.181) & K8s Worker 2 / Replica UPF pods  & 16\,GB & 256\,GB \\
core VM    (192.168.49.143) & Open5GS 5GC (AMF, SMF, NRF \ldots) & 8\,GB  & 128\,GB \\
UERANSIM VM (192.168.49.139)& gNB + UE Emulation               & 8\,GB  & 128\,GB \\
\midrule
\multicolumn{2}{l}{\textbf{Total}} & \textbf{64\,GB} & \textbf{1\,TB} \\
\bottomrule
\end{tabular}
\end{table}

\section{Software Requirements}

\begin{table}[H]
\centering
\caption{Software Stack and Version Requirements}
\label{tab:software}
\begin{tabular}{lll}
\toprule
\textbf{Component} & \textbf{Technology} & \textbf{Version} \\
\midrule
Operating System       & Ubuntu 22.04 LTS            & All nodes \\
Container Orchestrator & Kubernetes                  & v1.29.15 \\
Container Runtime      & containerd                  & v1.7.28 \\
CNI Plugin             & Flannel (VXLAN)             & v0.22 \\
5G Core                & Open5GS                     & v2.7.x \\
RAN / UE Emulation     & UERANSIM                    & v3.2.7 \\
LLM Backend            & Ollama                      & v0.3.x \\
LLM Model              & llama3.2:3b                 & 3B param, quantised \\
Agent Framework        & LangGraph                   & v0.1.x \\
Monitoring             & Prometheus                  & v2.45 \\
Visualisation          & Grafana                     & v10.x \\
Traffic Control        & Linux \texttt{tc} / HTB     & Kernel 6.8 \\
Programming Languages  & Python 3.10, Bash, YAML     & --- \\
\bottomrule
\end{tabular}
\end{table}
